\chapter{文献综述}


\section{探究性学习}


\subsection{国外相关研究}

18世纪法国杰出教育学家卢梭首次详细论述了发现法教学，重视纯实物教学，主张儿童通过观察自然界的原物获得真实感觉，再经过自主探究的方式获取学习，而不是通过直接告知或者间接教学，这是探究思想的前身。20世纪“现代教育学之父”杜威提出“五步教学法”，他认为教学阶段要遵从“情境—问题—观察—解决—实践”五步骤，这一理论适用于儿童从实践中手脑结合的探究学习。卢梭和杜威等人对发现教学以及探究教学的观点为后来探究性学习的提出与发展奠定了良好基础。

20世纪50年代，教育学家施瓦布指出，“如果要学生学习科学的方法，那么有什么学习比通过积极地投入到探究的过程中去更好呢？” 施瓦布对探究性学习和探究性教学作出解释，这句话大大推动了当时的探究性学习思潮。同时期后，布鲁纳创立了发现法，进而把发现法运用到美国教学中，取得了突出的成就。布鲁纳更注重探究性学习的理论依据，让教学更加科学和严谨。

施瓦布、布鲁纳等人的研究对各国50年代之后的课程及教材编写起到了重要影响，在课程学习中，发现学习、探究学习替代了原有的被动学习和阅读学习，学生真正参与到科学知识之中，通过问题和实践获得知识，更加强调学习科学的过程而非结果。

美国于1996年早期发布了《国家科学教育标准》，该文件把探索放在了科学研究的中心位置。科学探究有三方面含义，分别为学习内容、学习方式和教学指导思想\cite{科学探究与国家科学教育标准}。探究性学习在科学教育方面拥有正式文件的强力支持，与此同时，探究性学习成为了许多国家乐于选择的教育方式。


\subsection{国内相关研究}

探究性学习起源于国外，因此在其起源方面的国内研究甚少，直至陶行知等人将杜威的“教育即生活”等理论传播到中国，中国学者开始研究自由教育、探究教育等开放式教育。在此基础上，他提出“生活即教育”，生活与教育紧密结合，倡导在自然生活中进行教和学，重视自由探究的方法。
资料表明，我国于20世纪80年代就开始对探究学习与实验展开研究。美国L.Brenda教授举行有关探究研讨的教学讲学会并开展实验，提倡将探究学习作为小学教育的基础教育方式。在此之后，国内许多学者和教师都开始研究探究教育与探究实验，并有颇多成果。2001年底，新课改文件强调探究性学习应该成为被倡导的学习方法。张鹏(2003)用具体学校的实验案例进行分析，论述了探究性学习的具体实施策略和操作方法\cite{张鹏2003让探究性学习成为学生学习的基本方式}。他倡导让探究性学习成为小学学习的基本方法。彭建荣等(2007)讨论了探究性学习在网络环境下所形成的新的教学模式\cite{彭建荣2007论信息时代的探究性学习}。该文章结合信息时代，突出探究性学习在任何时代的适用性与融合性。高晓文等(2014)探究了教师的言行举止如何能够对自主探究课堂发挥积极作用\cite{高晓文2014教师的}。文章通过数据支持，说明了教师的“支持性”行为能够更好地促进探究性学习。该项研究意味着教师的教学方法和教学模式极其重要，即探究性学习不仅靠学生单方面自主思考、主动探究，辅佐以教师的支持与引导会更科学。近年来，21世纪技能被国际教育组织所强调，这是适应21世纪的核心基础教育能力，该技能的主要教学方法就是探究性学习。邓莉(2017)认为应当结合探究性学习和直接教学进行互补教学，探究性学习与21世纪技能息息相关\cite{邓莉2017如何在教学上落实}。



\section{数学探究性学习}

\subsection{国内相关研究}

2002年起，我国学者初步将探究性学习运用与数学教学中，从此我国学者开始重视探究性对于数学学习的作用。骆魁敏(2002)构建了高中数学网络探究性学习教学模式\cite{骆魁敏2003现代信息技术环境下高中数学研究性学习教学策略初探}。文章针对探究性教学模式提出操作特征并给出该模式的具体流程图，在探究性学习能够运用在数学教学上提供可能性。此后，更多学者认为探究性学习与数学教学的关系不可分割。由于数学学科的创新性和实践性，探究性学习是让学生主动参与获取相关数学知识，并不断提出新的问题与困惑，帮助学生更好地理解数学。随后，计算机技术运用到教学的大时代来临，如几何画板等各种数学软件被广泛应用于教学，探究性学习多与技术工具相结合，运用到中小学的数学教学中。王光明等(2013)利用Yves Chevallard所提出的“探索世界”范式对原本的探究性学习作出更进一步的思考，构建“Q-A-O”模式\cite{王光明2013探索世界}。文章提倡自主获取课外数学文献和建构命题及方法，并且强调学习数学的真正价值在于学生在探索的过程中感受数学的真实魅力。


\subsection{国外相关研究}

1970年开始，探究式教学的研究方兴未艾。许多美国教育家纷纷提出了实施探究性教学模式。之后，各国将探究作为基础课程的内容标准。但国外讨论的探究性学习大多用于问题解决和科学学习，与数学教学相关的文献较少，因此下文主要介绍探究性学习与数学学习的关系。Kim Taik H等(2010)评估了倡导探究式教学的STEP项目对学生学习数学的态度是否有影响\cite{kim2010impact}。他们通过实验发现探究性教学对学生的数学学习态度有积极的影响，尤其是先前对数学有抵触态度的学生受到的影响更为显著。Sulistiyo等(2020)探讨探究性学习对高中学生的计算思维和自我效能的影响\cite{sulistiyo2020effectiveness}。结果表明探究性学习对计算思维与自我效能感具有显著提升的效果，但在计算思维方面，探究性学习并没有成为比原有的更有效的学习模式，因为它在培养学生的提问能力上有欠缺，但在提高自我效能方面，由于探究性学习提供了更多样化的自我效能来源，因而比原先的学习模式更有效。Laudano等(2020)在数学算法研究中运用探究性学习方法。他们利用时长两年的教学实验说明探究性学习可以被教师应用到课堂中，为教学数学算法提供实用的策略和工具。并且该方法不仅能够对教学起到正向作用，还能够增强学习的趣味性，开阔视野\cite{laudano2020experience}。

随着新课改和新课标的要求，探究性学习被越来越多地用于高中数学教学中，众多学者在如何开展探究性学习方面作出了许多贡献，其中包含策略研究、实践研究、困惑与思考等等。但这些研究往往只宽泛地停留在高中数学教学的课堂上，并没有将课堂进行更细致的分类，或者仅仅强调新授课。因此，本人将研究高中数学习题课中的探究性学习，对数学课堂进行更细致的区分，结合习题课本身的特点与性质与探究性学习进行有效结合。



\section{高中数学课堂中的探究性学习}

笔者在本年3月底进入“中国知网”检索，以“探究性学习”为关键词，“高中数学”为主要主题，选择2012-2021年发表的文献，共检索到158篇文献，每年平均发表文献约15篇，其中大部分是期刊文献，少部分是会议文献与硕士论文。2012年的论文多关注高中数学中的探究性学习认知以及教学初探，2013-2016年进行教学策略初探即如何开展的问题，2017-2018年多进行教学设计实践以及困惑与思考，之后的论文大多关于探究性学习的教学实践与成效。以“探究性学习”为关键词，“数学习题课”为主要主题，共检索到6篇文献，这些文献都是期刊文献，其中近十年的有3篇，其主要内容大多为习题课的探究性学习实践。结合以上文献，有关高中数学教学的探究性学习的解析与应用已比较完备，从2012年的初探开始，每年的文献数量都保持在15篇以上，且2020年的数量达到31篇，这说明这个问题正处于持续研究阶段，但在不同的课型上研究还未完备，还有待广大学者进一步深究。